\section{ARIMA Model} 
\textit{Autoregressive Integrated Moving Average} (ARIMA) merupakan salah satu model deret waktu konvensional yang banyak digunakan, terutama sejak diperkenalkan oleh Box dan Jenkins (1970). 
Model ini membentuk pola deret waktu berdasarkan nilai masa lalu (lag) dan error acak (white noise)\citep{rathod2017forecasting}. 
ARIMA terdiri dari tiga komponen utama, yaitu: Autoregressive (AR), Integrated (I), dan Moving Average (MA). 
Komponen AR merepresentasikan pengaruh nilai sebelumnya, I berfungsi menjadikan data stasioner melalui proses differencing, sedangkan MA mencerminkan pengaruh kesalahan masa lalu. 
Secara umum, model ini dinotasikan sebagai ARIMA (p, d, q), di mana p adalah orde AR, d adalah tingkat differencing, dan q adalah orde MA.

Berikut ini adalah gambaran singkat mengenai model model dalam metode tersebut.
\begin{enumerate}
    \item Model Autoregressive (AR)
    \begin{equation}
        X_t = \phi_1 X_{t-1} + \phi_2 X_{t-2} + \cdots + \phi_p X_{t-p} + u_t
        \label{eq:ar_p_model}
    \end{equation}
    Model ini disebut model \textit{autoregressive} karena variabel $X_{t}$ bergantung secara linear pada nilai nya sendiri pada periode sebelumnya ($X_{t-1}$, $X_{t-2}$,..$X_{t-p}$). 
    Dengan kata lain, observasi saat ini $X_{t}$ dapat dinyatakan sebagai fungsi linear dari observasi sebelumnya berserta komponen error acak  yang dilambangkan dengan $U_{t}$
    \item Model Moving Averange (MA)
    \begin{equation}
        X_t = u_t - \theta_1 u_{t-1} - \theta_2 u_{t-2} - \cdots - \theta_q u_{t-q}
        \label{eq:ma_q_model}
    \end{equation}
    Model ini disebut model \textit{moving average} karena variabel $X_{t}$ bergantung secara linear pada error acak periode sekarang dan periode sebelumnya ($U_{t}$, $U_{t-1}$, $U_{t-2}$,...,$U_{t-q}$)
    \item Model Autoregressive Moving Average (ARMA)
    \begin{equation}
        X_t = \phi_1 X_{t-1} + \cdots + \phi_p X_{t-p} 
            - \theta_1 u_{t-1} - \cdots - \theta_q u_{t-q} + u_t
        \label{eq:arma_pq}
    \end{equation}
    Model ARMA menggabungkan kedua komponen yaitu ketergantungan terhadap nilai-nilai masa lalu ($X_{t-1}$, $X_{t-2}$,...,$X_{t-p}$) dan ketergantungan terhadap error sebelumnya ($U_{t}$, $U_{t-1}$,…,$U_{t-q}$).
    \item Model Autoregressive Integrated Moving Average (ARIMA)
    \begin{equation}
        \Phi(B)(1-B)^{d} W_{t} = \theta(B) U_{t}
        \label{eq:arma}
    \end{equation}
    dimana $W_{t}$ adalah diferensiasi pertama dari $X_{t}$, i.e., 
    $W_{t} = X_{t} - X_{t-1}$.
    Sebagian deret waktu bersifat tidak stasioner terhadap mean, sehingga diperlukan proses differensiasi untuk membuatnya stasioner. 
    Setelah proses differencing, deret yang telah stasioner dapat dimodelkan menggunakan ARMA, tetapi diperlukan operator differencing sebagai bagian dari model. 
    Ketika tingkat interasi pada sebuat proses menunjukkan berapa kali differencing dilakukan untuk mencapai kestasioneran, model tersebut dinamakana Autoregressive Integrated Moving Average (ARIMA) dengan bentuk umum ARIMA (p,d,q), dimana:
    \begin{itemize}
        \item p = orde autoregressive (AR)
        \item d = tingkat differencing
        \item q = orde moving average (MA)
    \end{itemize}
\end{enumerate}

Langkah langkah dalam membangunan model ARIMA yaitu:
\begin{enumerate}
    \item Estimasi Parameter Model: \\ 
    Latih model ARIMA dengan orde (p, d, q) yang telah ditentukan menggunakan data pelatihan. 
    Parameter-parameter model diestimasi menggunakan metode Maximum Likelihood Estimation (MLE).
    \item Menghitung Residual: \\
    Untuk memvalidasi model, residual dari persamaan yang diestimasi harus white noise. 
    Berikut adalah rumus ACF Residual
    \begin{equation}
        ACF(k) = 
        \frac{
            \displaystyle \sum_{t=1}^{T-k} (e_t - \bar{e})(e_{t+k} - \bar{e})
        }{
            \displaystyle \sum_{t=1}^{T-k} (e_t - \bar{e})^2
        }
        \label{eq:acf_formula}
    \end{equation}
    Keterangan:
    \begin{itemize}
        \item ACF(k) = Korelasi antara residu pada lag ke-k
        \item k = Lag yang diinginkan
        \item et = Residu pada waktu ke-t 
        \item e¯ = Rata-rata dari semua residu
        \item T = Jumlah observasi dalam deret waktu
    \end{itemize}
    \item Diagnostik Model Analisis Residual
    Periksa residu model (kesalahan prediksi) untuk memastikan mereka mendekati white noise (tidak ada autokorelasi, rata-rata nol, varians konstan). 
    Gunakan uji Ljung-Box dan plot ACF/PACF residu. Rumus Ljung-Box :
    \begin{equation}
        Q = n(n+2) \sum_{k=1}^{h} \frac{\hat{\rho}_k^{\,2}}{n - k}
        \label{eq:ljung_box_q}
    \end{equation}
    Keterangan:
    \begin{itemize}
        \item Q: Statistik Uji Ljung-box
        \item n : Jumlah observasi (ukuran sampel) dalam deret waktu.
        \item h : Jumlah lag yang diuji (batas atas jumlah autokorelasi yang dihitung). Nilai h dapat ditentukan sesuai kebutuhan; secara umum digunakan nilai h = 1,2,3
        \item $\rho$k: Estimasi koefisien autokorelasi dari residu pada lag ke-k
    \end{itemize}
\end{enumerate}